\vspace{5pt plus 4pt minus 1pt}
\noindent{\renewcommand{\arraystretch}{1.24}\setlength{\tabcolsep}{4pt}\rowcolors{1}{kbvrow}{white}
\begin{tabularx}{\textwidth}{X X X}
ತೃ. ವಿ. & ಅಗಸನಿಂದ & ಅಗಸರಿಂದ \\
ಚ. ವಿ. & ಅಗಸನಿಗೆ & ಅಗಸರಿಗೆ \\
ಪಂ. ವಿ. & ಅಗಸನ ದೆಸೆಯಿಂದ & ಅಗಸರ ದೆಸೆಯಿಂದ \\
ಷ. ವಿ. & ಅಗಸನ & ಅಗಸರ \\
ಸ. ವಿ. & ಅಗಸನಲ್ಲಿ & ಅಗಸರಲ್ಲಿ \\
ಸಂ. ವಿ. & ಅಗಸನೇ ! & ಅಗಸರೇ ! \\
\end{tabularx}}
\vspace{5pt plus 4pt minus 1pt}

\kbvsechead{ಇಕಾರಾಂತ ಸ್ತ್ರೀಲಿಂಗ ಶಬ್ದ “ರಾಣಿ” :}

\vspace{5pt plus 4pt minus 1pt}
\noindent{\renewcommand{\arraystretch}{1.24}\setlength{\tabcolsep}{4pt}\rowcolors{2}{kbvrow}{white}
\begin{tabularx}{\textwidth}{X X X}
{\kbvsecfont\color{kbvaccentdark}ವಿ. ಪ್ರ.} & {\kbvsecfont\color{kbvaccentdark}ಏ. ವ.} & {\kbvsecfont\color{kbvaccentdark}ಬ. ವ.} \\[1pt]
\arrayrulecolor{kbvline}\specialrule{0.5pt}{0pt}{2pt}\arrayrulecolor{black}
ಪ್ರ. ವಿ. & ರಾಣಿಯು & ರಾಣಿಯರು \\
ದ್ವಿ. ವಿ. & ರಾಣಿಯನ್ನು & ರಾಣಿಯರನ್ನು \\
ತೃ. ವಿ. & ರಾಣಿಯಿಂದ & ರಾಣಿಯರಿಂದ \\
ಚ. ವಿ. & ರಾಣಿಗೆ & ರಾಣಿಯರಿಗೆ \\
ಪಂ. ವಿ. & ರಾಣಿಯ ದೆಸೆಯಿಂದ & ರಾಣಿಯರ ದೆಸೆಯಿಂದ \\
ಷ. ವಿ. & ರಾಣಿಯ & ರಾಣಿಯರ \\
ಸ. ವಿ. & ರಾಣಿಯಲ್ಲಿ & ರಾಣಿಯರಲ್ಲಿ \\
ಸಂ. ವಿ. & ರಾಣಿಯೇ ! & ರಾಣಿಗಳೇ ! \\
\end{tabularx}}
\vspace{5pt plus 4pt minus 1pt}

\kbvsechead{ಉಕಾರಾಂತ ನಪುಂಸಕ ಲಿಂಗ ಶಬ್ದ “ದೇವರು” :}

\vspace{5pt plus 4pt minus 1pt}
\noindent{\renewcommand{\arraystretch}{1.24}\setlength{\tabcolsep}{4pt}\rowcolors{2}{kbvrow}{white}
\begin{tabularx}{\textwidth}{X X X}
{\kbvsecfont\color{kbvaccentdark}ವಿ. ಪ್ರ.} & {\kbvsecfont\color{kbvaccentdark}ಏ. ವ.} & {\kbvsecfont\color{kbvaccentdark}ಬ. ವ.} \\[1pt]
\arrayrulecolor{kbvline}\specialrule{0.5pt}{0pt}{2pt}\arrayrulecolor{black}
ಪ್ರ. ವಿ. & ದೇವರು & ದೇವರುಗಳು \\
ದ್ವಿ. ವಿ. & ದೇವರನ್ನು & ದೇವರುಗಳನ್ನು \\
ತೃ. ವಿ. & ದೇವರಿಂದ & ದೇವರುಗಳಿಂದ \\
ಚ. ವಿ. & ದೇವರಿಗೆ & ದೇವರುಗಳಿಗೆ \\
ಪಂ. ವಿ. & ದೇವರ ದೆಸೆಯಿಂದ & ದೇವರುಗಳ ದೆಸೆಯಿಂದ \\
ಷ. ವಿ. & ದೇವರ & ದೇವರುಗಳ \\
ಸ. ವಿ. & ದೇವರಲ್ಲಿ & ದೇವರುಗಳಲ್ಲಿ \\
ಸಂ. ವಿ. & ದೇವರೇ ! & ದೇವರುಗಳೇ ! \\
\end{tabularx}}
\vspace{5pt plus 4pt minus 1pt}
